\section{Proof of Theorem~\ref{thm:bounded}}\label{sec:thmA}
In this section, $G=K_n$ with $n$ large, $W=I$ (so that $H_A=L_A$), and
\[
B=2I-\tfrac2nJ,\qquad\alpha=\tfrac1{100},\qquad\beta=\tfrac25,\qquad\ell=\lceil6250\log n\rceil,\qquad q_0=\lceil10^9\ell\log n\rceil .
\]
We apply Lemmas~\ref{lem:sample} and~\ref{lem:extract} with $k=n-1$ and $m=\lfloor n/2\rfloor$, so that the forests we find are spanning trees.
Note that, for even $n$, $B=L_{K_n}/(n/2)$, and that $\ell=\lceil1000\beta^{-2}\log n\rceil$ and $\alpha^2q_0\ge10^5\ell\log n$.

The key point is that closeness of $L_R/q$ to $B$ gives all the conditions of Lemma~\ref{lem:sample}.
For small sets, it gives an upper bound on the number of edges inside them, and, for large sets, an upper bound via the number of edges meeting their complements.
The idea of deducing conditions for packing spanning trees from eigenvalues goes back to Cioab\u{a} and Wong~\cite{CW}.

\begin{lemma}\label{lem:ranks}
Let $(R,q)$ be good, as in Definition~\ref{def:good}. Then the following hold.
\begin{props}
\item $(2-\frac2n-\alpha)q\le d_R(v)\le(2+\alpha)q$ for every vertex $v$.\label{rk:deg}
\item $e_R(X)\le\frac35q(|X|-1)$ for every $X$ with $2\le|X|\le n/2$.\label{rk:small}
\item $I_R(U)\ge\frac75q|U|$ for every $U$ with $1\le|U|\le n/2$.\label{rk:large}
\end{props}
\end{lemma}
\begin{proof}
Let $A=L_R/q-B$, so that $\|A\|\le\alpha$. Part~\ref{rk:deg} follows from $d_R(v)/q=(L_R/q)_{vv}=2-\frac2n+A_{vv}$.
For~\ref{rk:small}, let $|X|=s$. As $\one_X^{\!\top}B\one_X=2s-2s^2/n$, \eqref{eq:edgecount} gives
\[
\frac{2e_R(X)}q=\frac{2s(s-1)}n+\sum_{v\in X}A_{vv}-\one_X^{\!\top}A\one_X\le\frac{2s(s-1)}n+2\alpha s .
\]
If $4\le s\le n/2$, then $s(s-1)/n\le(s-1)/2$ and $s\le\frac43(s-1)$, so $e_R(X)/q<\frac35(s-1)$. If $s\in\{2,3\}$, then $e_R(X)\le3\le\frac35q(s-1)$ as $q\ge q_0$.
For~\ref{rk:large}, let $|U|=u\le n/2$. By~\ref{rk:deg} and the bound $e_R(U)/q\le u(u-1)/n+\alpha u$ just shown,
\[
\frac{I_R(U)}q=\sum_{v\in U}\frac{d_R(v)}q-\frac{e_R(U)}q\ge u\Bigl(2-\frac{u+1}n-2\alpha\Bigr)>\frac75u . \qedhere
\]
\end{proof}

\begin{proof}[Proof of Theorem~\ref{thm:bounded}]
We first check that every good pair $(R,q)$ with $q>q_0$ satisfies the conditions of Lemma~\ref{lem:sample}.
Condition~\ref{S:count} is part of being good, and~\ref{S:deg} holds by Lemma~\ref{lem:ranks}\ref{rk:deg}.
For~\ref{S:rank}, recall that $k=n-1$, and let $2\le|X|\le n-1$. If $|X|\le n/2$, then $e_R(X)\le\frac35q(|X|-1)$ by Lemma~\ref{lem:ranks}\ref{rk:small}. Otherwise $U=V\setminus X$ has $|U|<n/2$, and $I_R(U)\ge\frac75q|U|$ by Lemma~\ref{lem:ranks}\ref{rk:large}.
Thus the hypothesis of Lemma~\ref{lem:extract} holds.

If $n$ is even, let $R_m=E(K_n)$, so that $L_{R_m}/m=B$.
If $n=2m+1$ is odd, let $R_m=E(K_n)\setminus M$ for a matching $M$ with $m$ edges. Then $|R_m|=m(n-1)$, and every colour has at most $\lfloor n/2\rfloor=m$ edges. Moreover, $L_{K_n}/m-B$ vanishes on $\one$ and equals $\frac1mI$ on $\one^\perp$, and $\|L_M\|\le2$, so $\|L_{R_m}/m-B\|\le3/m\le\alpha/4$.
In both cases, Lemma~\ref{lem:extract} gives $m-q_0=\lfloor n/2\rfloor-O(\log^2n)$ edge-disjoint rainbow spanning trees.
\end{proof}

For even $n$, Lemma~\ref{lem:extract} also tells us that the remaining edge set $R$ is such that $(R,q_0)$ is good, so $\Delta(R)\le3q_0$ and every colour has at most $q_0$ edges in $R$.
Each edge of $R$ therefore shares a vertex or a colour with fewer than $7q_0$ other edges of $R$, and colouring greedily splits $R$ into at most $7q_0$ rainbow matchings.
Together with the trees, this partitions $E(K_n)$ into $n/2+O(\log^2n)$ rainbow forests, as claimed in Section~\ref{sec:bounded}.
